\section{General linear Lie algebras}
\label{section:general-linear-algebras}
In this section, we study the Lie algebra of general linear groups. We shall prove that, in fact, such a Lie algebra $\gl_\A$ coincides with $\lnr_\A^-$, which is the differential Lie algebra of the maximal subring of the LEC object $\lnr_\A$, whose Lie bracket is given by the commutator. In particular, when $\A$ has negatives, this makes $\gl_\A$ isomorphic to $\lnr_\A$ itself, recapturing the classical result of differential geometry.



%__________________________________________________________________________
\subsection{The Lie algebra of a general linear group object}
\label{subsection:general-linear-algebra}
Recall that, as explained in the introduction, a \textbf{Lie group} in a Cartesian tangent category consists of a group object $G$ which admits the tangent space at the unit, in that the tangent pullback $\T_\e G$ of the tangent bundle $p_G\colon\T G\to G$ along the unit map $\e\colon\*\to G$ exists~\cite[Definition~2.13]{lanfranchi:lie-correspondence-tangent-cats}.

Recall also that the Lie functor $\g(-)$ of $G$ is the functor that sends each $M$ to the Lie algebra $\g(M)$ of $M$-parametrized left-invariant vector fields ($M$-\PLIVFs for short) of $G$. The Lie algebra of a Lie group is the representing object $\g$ of the Lie functor $\g(-)$ of $G$, which coincides with the tangent space $\T_\e G$. By~\cite[Proposition~4.3]{lanfranchi:lie-correspondence-tangent-cats}, this mounts to say that there is a \PLIVF $\omega\colon G\times\g\to\T G$ which is universal, in that for any $M$-\PLIVF $X\colon G\times M\to\T G$ of $G$ there exists a unique morphism $\hat X\colon M\to\g$ such that $X=(\id\times\hat X)\omega$. The universal \PLIVF of a Lie group is determined by the inclusion $\iota_\e\colon\T_\e G\to\T G$.

\begin{lemma}[{\cite[Lemma~4.5]{lanfranchi:lie-correspondence-tangent-cats}}]
\label{lemma:omega-vs-iota}
For a Lie group $G$, the following identity holds:
\begin{align*}
&\omega_G=(z\times\iota_\e)\T\m_G
\end{align*}
where $\omega_G\colon G\times\T_\e G\to\T G$ denotes the universal \PLIVF of $G$ and $\iota_\e\colon\T_\e G\to\T G$ is the inclusion of the tangent space in the tangent bundle of $G$.
\end{lemma}

Our goal is to compare the differential Lie algebra $\gl_\A$ of the general linear group $\GL_\A$ with the differential Lie algebra $\lnr_\A^-$ of Proposition~\ref{proposition:LEC-lie-algebra}. To make this comparison, we shall first re-interpret the existence of the Lie algebra object $\lnr_\A^-$ via the representability of the functor
\begin{align*}
&\lnr_\A^-(-)\colon\X\to\Set^\op
\end{align*}
that sends each object $M$ to the set $\lnr_\A^-(M)$ of $M$-parametrized linear endomorphisms $f\colon\A\times M\to\A$ of $\A$ for which there exists a morphism $-f$ such that, their sum $f+(-f)\=\<f,-f\>+$ vanishes, that is, $f+(-f)=!0$.

Another approach to construct the functor $\lnr_\A^-(-)$ consists of considering first the functor
\begin{align*}
&\lnr_\A(-)\colon\X\to\Set^\op  &&M\mapsto\lnr_\A(M)\=\DO_\multi(\Simple_M[\X,\TT])(\A,\A)
\end{align*}
that sends each $M$ to the set of $M$-parametrized linear endomorphisms. When the LEC object of $\A$ exists, this functor is representable and represented by the LEC object $\lnr_\A$. Moreover, since $M$-parametrized linear endomorphisms turn out to be simply linear endomorphisms in $\Simple_M[\X,\TT]$, each set $\lnr_\A(M)$ comes equipped with a rig structure that turns the functor $\lnr_\A(-)$ into a rig object in the presheaf category of $\X$. Since the presheaf category is always complete, one can construct the maximal subring of $\lnr_\A(-)$ in the presheaf category by taking the equalizer of the sum $+$ with the zero map $!0$ of $\lnr_\A(-)$. Thus, $\lnr_\A^-(-)$ is precisely the resulting ring object in the presheaf category.\\

The Lie functor $\gl(-)$ of the general linear group $\GL_\A$ of a differential object $\A$ is related to the functor $\lnr_\A(-)$ via a natural Lie algebra homomorphism, as proved by the following proposition.

\begin{proposition}
\label{proposition:the-map-chi}
Suppose that $\A$ is a differential object $\A$ which admits the general linear group, and that the tangent bundle $\T\GL_\A$ admits negatives\footnote{Negatives are required to define the Lie functor a group object. Notice that the tangent bundle of a Lie group has always negatives.}. Therefore, there is a Lie algebra homomorphism
\begin{align*}
&\chi_M^\A\colon\gl_\A(M)\to\lnr_\A^-(M)
\end{align*}
natural in $M$ which sends each $M$-\PLIVF $X\colon\GL_\A\times M\to\T\GL_\A$ to:
\begin{align*}
&\varphi_X\colon \A\times M\xrightarrow{\id\times\e\times\id}\A\times\GL_\A\times M\xrightarrow{z\times X}\T(\A\times\GL_\A)\xrightarrow{\T\Ev}\T\A\xrightarrow{\hat p}\A
\end{align*}
\end{proposition}
\begin{proof}
For starters, we show that $\varphi_X$ is an $M$-parametrized linear endomorphism of $\A$:
\begin{align*}
\varphi_X\lambda_\A&=~(\id\times\e\times\id)(z\times X)\T\Ev\hat p\lambda_\A    \tag{Definition of $\varphi_X$}\\
&=~(\id\times\e\times\id)(z\times X)\T\Ev\T\lambda_\A c\T\hat p         \tag{$\hat p\lambda_\A=\T\lambda_\A c\T\hat p$}\\
&=~(\id\times\e\times\id)(z\T\lambda_\A\times X\T z)\T\T\Ev c\T\hat p     \tag{Linearity of $\Ev$}\\
&=~(\id\times\e\times\id)(\lambda_\A z\times z\T X)\T\T\Ev\T\hat p      \tag{Naturality of $c$ and $z$}\\
&=~(\lambda_\A\times z)\T((\id\times\e\times\id)(z\times X)\T\Ev)\T\hat p\\
&=~(\lambda_\A\times z)\T\varphi_X
\end{align*}
Thus, we constructed a function $\chi_M^\A\colon\gl_\A(M)\to\lnr_\A(M)$ from the Lie functor to $\lnr_\A(M)$, which is the set of $M$-parametrized linear endomorphisms of $\A$. However, since, by assumption, the tangent bundle of $\GL_\A$ admits negatives, such a function factors through the natural inclusion $\lnr_\A^-(M)\to\lnr_\A(M)$. We shall now prove that this function preserves the Lie algebra structure. We start by showing that $\varphi_0$ is the zero morphism, where $0\=\pi_1z\colon\GL_\A\times M\to\T\GL_\A$:
\begin{align*}
\varphi_0&=~(\id\times\e\times\id)(z\times(\pi_1z)\T\Ev\hat p  \tag{Definition of $\varphi_0$}\\
&=~\pi_1(\id\times\e)z\T\Ev\hat p                  \tag{Naturality of $\pi_1$}\\
&=~\pi_1(\id\times\e)\Ev z\hat p                                 \tag{Naturality of $z$}\\
&=~\pi_1(\id\times\e)\Ev!0                                      \tag{$z\hat p=!0$}\\
&=~!_{\id\times M}0\\
&=~0
\end{align*}
Now, let us show that $\chi$ preserves the sum. Given two \PLIVFs $X,Y\colon\GL_\A\times M\to\T\GL_\A$:
\begin{align*}
\varphi_{X+Y}&=~(\id\times\e\times\id)(z\times(\<X,Y\>s))\T\Ev\hat p      \tag{Definition of $\varphi_{X+Y}$}\\      
&=~(\id\times\e\times\id)(\<z\times X,z\times Y\>)\T_2\Ev s\hat p  \tag{Linearity of $\Ev$ implies additivity of $\Ev$}\\
&=~(\id\times\e\times\id)(\<z\times X,z\times Y\>)\T_2\Ev\T_2\hat p+  \tag{$s\hat p=\T_2\hat p+$}\\
&=~\left\<(\id\times\e\times\id)(z\times X)\T\Ev\T\hat p,(\id\times\e\times\id)(z\times Y)\T\Ev\T\hat p\right\>+\\
&=~\varphi_X+\varphi_Y
\end{align*}
Finally, we shall prove that $\varphi_{[X,Y]}=[\varphi_X,\varphi_Y]$. Let us begin with considering the composition $\varphi_Y;_\KL\varphi_X$, where composition is in the coKleisli category $\Simple_M(\X,\TT)$:
\begin{align*}
\varphi_Y;_\KL\varphi_X&(\id\times\Delta_M)(\varphi_Y\times\id)\varphi_X
\end{align*}
After some basic rewriting and using the naturality of $z$, we can rewrite this as:
\begin{align*}
&\left((z\times\T\e\times z)(\T z\times\T Y)\T\T\Ev\T\hat p\times(\e\times\id)X\right)\T\Ev\hat p
\end{align*}
However, since $\Ev$ is linear in the $\A$-entry, $(\hat p\times\id)\Ev=(\id\times z)\T\Ev\hat p$. Using this identity, we can rewrite $\varphi_Y;_\KL\varphi_X$ as follows:
\begin{align*}
&\left((z\times\T\e\times z)(\T z\times\T Y)\T\T\Ev\times(\e\times\id)X\T z\right)\T\T\Ev\T\hat p\hat p
\end{align*}
However, using the left-invariance of $Y$, we can write the following commutative diagram:
\begin{equation*}
% https://q.uiver.app/#q=WzAsNixbMCwwLCJNIl0sWzAsMSwiXFxHTF9BXFx0aW1lcyBNXFx0aW1lc1xcR0xfQVxcdGltZXMgTVxcdGltZXMgTSJdLFswLDIsIlxcVChcXEdMX0FcXHRpbWVzXFxHTF9BXFx0aW1lcyBNKSJdLFswLDMsIlxcVFxcVChcXEdMX0FcXHRpbWVzXFxHTF9BKSJdLFsxLDIsIlxcVChcXEdMX0FcXHRpbWVzIE0pIl0sWzEsMywiXFxUXFxUXFxHTF9BIl0sWzAsMSwiXFw8XFxlLFxcaWRfTSxcXGUsXFxpZF9NLFxcaWRfTVxcPiIsMl0sWzEsMiwiWFxcaV97XFxHTF9BfVxcdGltZXMgWFxcdGltZXMgel9NIiwyXSxbMiwzLCJcXFQoel97XFxHTF9BfVxcdGltZXMgWSkiLDJdLFsyLDQsIlxcVChcXG1fe1xcR0xfQX1cXHRpbWVzXFxpZF9NKSJdLFs0LDUsIlxcVCBZIl0sWzMsNSwiXFxUXFxUXFxtX3tcXEdMX0F9IiwyXSxbMCw0LCJcXDxcXFRcXGUsel9NXFw+IiwwLHsiY3VydmUiOi00fV1d
\begin{tikzcd}
M & \\
{\GL_\A\times M\times\GL_\A\times M\times M} \\
{\T(\GL_\A\times\GL_\A\times M)} & {\T(\GL_\A\times M)} \\
{\T\T(\GL_\A\times\GL_\A)} & {\T\T\GL_\A}
\arrow["{\<\e,\id,\e,\id,\id\>}"', from=1-1, to=2-1]
\arrow["{\<\T\e,z\>}", curve={height=-24pt}, from=1-1, to=3-2]
\arrow["{X\i\times X\times z}"', from=2-1, to=3-1]
\arrow["{\T(\m\times\id)}", from=3-1, to=3-2]
\arrow["{\T(z\times Y)}"', from=3-1, to=4-1]
\arrow["{\T Y}", from=3-2, to=4-2]
\arrow["{\T\T\m}"', from=4-1, to=4-2]
\end{tikzcd}
\end{equation*}
Furthermore, by construction, the multiplication $\m$ of $\GL_\A$ is uniquely defined by the following identity, $(\id\times\m)\Ev=(\Ev\times\id)\Ev$. Using this identity and the left-invariance of $Y$, we can rewrite $\varphi_Y;_\KL \varphi_X$ as follows:
\begin{align*}
&(z\times\e\times\id\times\e\times\id\times\e\times\id)(z\times((X\T z\T\i)\times((X\times z)\T Y)\times(X\T z)))\\
&\qquad(\T z\times\T\T\m_3)\T\T\Ev\T\hat p\hat p
\end{align*}
where $\m_3\=(\id\times\m)\m=(\m\times\id)\m$. In this expression there are three terms, $(X\T z\T\i)$, $((X\times z)\T Y)$, and $(X\T z)$, multiplied together by $\T\T\m_3$. However, the first and last terms invert each other. Furthermore, by an Eckmann-Hilton argument, the derivative of the multiplication $\m$ coincides with the sum of tangent vectors and thus is commutative. This allows us to swap the terms and simplify the former with the latter. What is left is the following expression:
\begin{align*}
&\varphi_Y;_\KL \varphi_X=((\id\times\e\times\id)(z\times X)\times z)(\T z\times\T Y)\T\T\Ev\T\hat p\hat p\\
&\qquad=
(\id\times\e\times\id)(z\T z\times(X;_M\T_MY))\T\T\Ev\T\hat p\hat p
\end{align*}
where $X;_M\T_MY$ denotes the composition of $X\colon\GL_\A\nto\T\GL_\A$ with $\T_MY\colon\T\GL_\A\nto\T\T\GL_\A$ in $\Simple_M(\X,\TT)$. Similarly, we also obtain that:
\begin{align*}
&\varphi_X;_\KL \varphi_Y=(\id\times\e\times\id)(z\T z\times(Y;_\KL\T^\KL X))\T\T\Ev\T\hat p\hat p
\end{align*}
Since $\A$ is a differential object, $\T\T\A\cong\A\times\A\times\A\times\A$. Thus, the canonical flip $c$ swaps the internal components, while $\T\hat p\hat p\colon\T\T\A\to\A$ projects the fourth component down to $\A$. Therefore, $\T\hat p\hat p=c\T\hat p\hat p$ (\cite[Proposition~3.6]{cockett:differential-bundles}). So, we can also rewrite:
\begin{align*}
&\varphi_X;_\KL \varphi_Y=(\id\times\e\times\id)(z\T z\times(Y;_\KL\T^\KL X)c)\T\T\Ev\T\hat p\hat p
\end{align*}
where we used the naturality of $c$ and that $z\T z c=z z=z\T z$. Now, let us unwrap the definition of $\varphi_{[X,Y]}$:
\begin{align*}
\varphi_{[X,Y]}&=~(\id\times\e\times\id)(z\times[X,Y])\T\Ev\hat p\\
&=~(\id\times\e\times\id)\left\{\left(\left(z\T z\times(X;_\KL\T^\KL Y)\right)-\left(z\T z\times(Y;_\KL\T^\KL Xc)\right)\right)\T\T\Ev\right\}\hat p
\end{align*}
By~\cite[Example~3.10]{cockett:connections}, if $\A$ is a differential object, $\<p_{\T \A}p_\A,T\hat p\hat p\>\colon\T\T\A\to\T\A$ defines a vertical affine connection of $\A$\footnote{Notice that in~\cite[Example~3.10]{cockett:connections}, the first component of $\T\A\cong\A\times\A$ contain the tangent vectors and not the points. We use the opposite convention.}. However, by~\cite[Lemma~3.13]{cockett:connections}, if $K\colon\T\T M\to\T M$ is an affine connecection, then $\{\varphi\}=\varphi K$. However, we can also compute:
\begin{align*}
&\{\varphi\}\hat p=\varphi\<p_{\T\A}p_\A,\T\hat p\hat p\>\hat p=\varphi\T\hat p\hat p
\end{align*}
Therefore, we can rewrite $\varphi_{[X,Y]}$ as:
\begin{align*}
&\left((\id\times\e\times\id)\left(z\T z\times(X;_\KL\T^\KL Y)\right)-(\id\times\e\times\id)\left(z\T z\times(Y;_\KL\T^\KL Xc)\right)\right)\\
&\qquad\T\T\Ev\T\hat p\hat p
\end{align*}
Finally, using the additivity of $\hat p$, we conclude that $\varphi_{[X,Y]}=\varphi_Y\varphi_X-\varphi_X\varphi_Y=[\varphi_X,\varphi_Y]$.
\end{proof}

The Lie algebra morphism $\chi\colon\gl_\A(M)\to\lnr_\A^-(M)$ compares the Lie functor of the general linear group $\GL_\A$ with the functor $\lnr_\A^-(-)$ that classifies linear endomorphisms of $\A$. When $\A$ admits the LEC object and its maximal subring $\lnr_\A^-$, then the functor $\lnr_\A^-(-)$ becomes representable. Furthermore, if also the general linear group $\GL_\A$ of $\A$ is a Lie group, then its Lie functor $\gl_\A(-)$ is also representable. We shall introduce a useful definition.

\begin{definition}
\label{definition:general-linear-lie-group}
We shall say that a differential object $\A$ admits the \textbf{general linear Lie group} provided that the general linear group object $\GL_\A$ of $\A$ exists and is a Lie group. We shall denote the differential Lie algebra of $\GL_\A$ by $\gl_\A$.
\end{definition}

We shall now prove that the comparison natural transformation $\chi$ is completely determined by the inclusion map $\theta_\A\colon\GL_\A\to\lnr_\A$, uniquely defined by $(\id\times\theta)\ev=\Ev$. To this end, we shall denote by $\rho\colon\lnr_\A^-\to\lnr_\A$ the inclusion map. Recall also that $\gl_\A\cong\T_\e\GL_\A$ and that $\lnr_\A\cong\T_\e\lnr_\A$, since $\lnr_\A$ is a differential object.

\begin{lemma}
\label{lemma:chi-vs-theta-maps}
If $\A$ admits the LEC object $\lnr_\A$, its maximal subring $\lnr_\A^-$, and the general linear Lie group $\GL_\A$, then there exists a linear Lie algebra homomorphism
\begin{align*}
&\chi_\A\colon\gl_\A\to\lnr_\A^-
\end{align*}
induced by the natural transformation $\chi_\A^M\colon\gl_\A(M)\to\lnr_\A(M)$. Furthermore:
\begin{align*}
&\chi_\A\rho_\A=\T_\e\theta_\A
\end{align*}
\end{lemma}
\begin{proof}
By the Yoneda lemma, there is a unique morphism $\chi'_\A\colon\gl_\A\to\lnr_\A$, which renders the following diagram
\begin{equation*}
% https://q.uiver.app/#q=WzAsMyxbMCwwLCJcXEFcXHRpbWVzXFxnbF9cXEEiXSxbMCwxLCJcXEFcXHRpbWVzXFxsbnJfXFxBIl0sWzEsMSwiXFxBIl0sWzEsMiwiXFxldl9cXEEiLDJdLFswLDEsIlxcaWRfXFxBXFx0aW1lc1xcY2hpX1xcQSciLDJdLFswLDIsIlxcdmFycGhpX1xcb21lZ2EiXV0=
\begin{tikzcd}
{\A\times\gl_\A} & \\
{\A\times\lnr_\A} & \A
\arrow["{\id\times\chi_\A'}"', from=1-1, to=2-1]
\arrow["{\varphi_\omega}", from=1-1, to=2-2]
\arrow["{\ev}"', from=2-1, to=2-2]
\end{tikzcd}
\end{equation*}
commutative, where $\varphi_\omega=\chi_{\gl_\A}^\A(\omega)$ and where $\omega\colon\GL_\A\times\gl_\A\to\T\GL_\A$ is the universal \PLIVF of $\GL_\A$. Let us unwrap the definition of $\varphi_\omega$:
\begin{align*}
\varphi_\omega&=~(\id\times\e\times\id)(z\times\omega)\T\Ev\hat p       \tag{Definition of $\varphi_\omega$}\\
&=~(\id\times\e\times\id)(z\times(z\times\iota_\e)\T\m)\T\Ev\hat p   \tag{Lemma~\ref{lemma:omega-vs-iota}}\\
&=~(\id\times\e\times\id)((z\times z)\T\Ev\times\iota_\e)\T\Ev\hat p   \tag{Definition of $\m$}\\
&=~((\id\times\e)\Ev z\times\iota_\e)\T\Ev\hat p      \tag{Naturality of $z$}\\
&=~(z\times\iota_\e)\T\Ev\hat p                          \tag{Definition of $\e$}
\end{align*}
Now, let us consider $\T_\e\theta_\A$, which is defined as the unique morphism, such that $\T_\e\theta_\A\<!\e,\lambda_{\lnr_\A}\>=\iota_\e\T\theta_\A$. Thus, by postcomposing on both sides by $\hat p_{\lnr_\A}$, we obtain that, $\T_\e\theta_\A=\T_\e\theta_\A\<!\e,\lambda_{\lnr_\A}\>\hat p_{\lnr_\A}=\iota_\e\T\theta_\A\hat p_{\lnr_\A}$, where we used that $\lambda_{\lnr_\A}\hat p_{\lnr_\A}$ is the identity. Thus:
\begin{align*}
(\id\times\T_\e\theta_\A)\ev&=~(\id\times(\iota_\e\T\theta_\A\hat p_{\lnr_\A}))\ev\\
&=~(\id\times\iota_\e\T\theta_\A)(\id\times\hat p_{\lnr_\A})\ev\\
&=~(\id\times\iota_\e\T\theta_\A)(z\times\id)\T\ev\hat p     \tag{Bilinearity}\\
&=~(z\times\iota_\e)\T((\theta_\A\times\id)\ev)\hat p\\
&=~(z\times\iota_\e)\T\Ev\hat p                                          \tag{Definition of $\theta_\A$: $(\theta_\A\times\id)\ev=\Ev$}
\end{align*}
Using the universal property of $\ev$, we conclude that $\chi_\A'=\T_\e\theta_\A$. However, since every morphism $\T_\e f$ is always linear, $\chi_\A'$ is also linear. To conclude, it is only left to show that $\chi_\A'=\chi_\A\rho$ for some $\chi_\A\colon\gl_\A\to\lnr_\A^-$. However, since $\gl_\A$ has negatives given by $\T_\e\i$, where $\i$ denotes the inverse map of $\GL_\A$, it follows that $\chi_\A'$ lands inside of $\lnr_\A^-$, that is, it factors through $\rho\colon\lnr_\A^-\to\lnr_\A$. Using the universal property of $\rho$ that establishes that $\<\rho,-\rho\>$ is an equalizer map of $+$ and $!0$ as in Definition~\ref{definition:maximal-differential-subring}, one can easily prove that $\chi_\A$ is in fact a linear Lie algebra homomorphism.
\end{proof}

Before proving the main result, we indulge in another technical lemma.

\begin{lemma}
\label{lemma:multiplication-and-unit-at-e}
By identifying the tangent space $\T_\e\lnr_\A$ of $\lnr_\A$ at $\e\colon\*\to\lnr_\A$ with $\lnr_\A$ itself, the morphism $\T_\e\m$ coincides with the sum $+$ of $\lnr_\A$ and $\T_\e\e$ with the zero $0$.
\end{lemma}
\begin{proof}
To prove this lemma, one can argue as in~\cite[Theorem~2]{cockett:lie-groups-tangent-cats}, in which Cockett and Schwarz used an Eckmann-Hilton argument to prove that $\T_\e\m=+$ and $\T_\e\e=0$, for Lie groups.
\end{proof}

Thanks to Theorem~\ref{theorem:triple-equalizer}, the general linear group $\GL_\A$ is the tangent triple equalizer of $\m$, $!\e$, and $\m^\op$ of $\lnr_\A$. Using that result, we now prove the main theorem of this paper: that the differential Lie algebra of $\GL_\A$ is isomorphic to the differential Lie algebra $\lnr_\A^-$.

\begin{theorem}
\label{theorem:main-theorem}
Suppose that $\A$ admits the general linear Lie group object $\GL_\A$ and the LEC object $\lnr_\A$. Then, the maximal subring $\lnr_\A^-$ of $\lnr_\A$ exists, and the differential Lie algebra $\gl_\A$ of $\GL_\A$ is isomorphic to the differential Lie algebra $\lnr_\A^-$. In particular, the map $\chi$ of Lemma~\ref{lemma:chi-vs-theta-maps} is an isomorphism of differential Lie algebras.
\end{theorem}
\begin{proof}
From Theorem~\ref{theorem:triple-equalizer}, $\GL_\A$ can be realized as the tangent triple equalizer of $\m$, $!\e$, and $\m^\op$ of $\lnr_\A$. By applying $\T_\e$, which takes the tangent space at $\e$ to this tangent triple equalizer, we obtain a new tangent triple equalizer:
\begin{equation*}
% https://q.uiver.app/#q=WzAsNCxbMCwwLCJcXFRfXFxlXFxHTF9cXEEiXSxbMSwwLCJcXFRfXFxlXFxsbnJfXFxBXFx0aW1lc1xcVF9cXGVcXGxucl9cXEEiXSxbMywwLCJcXFRfXFxlXFxsbnJfXFxBIl0sWzIsMCwiXFwqIl0sWzAsMV0sWzEsMiwiXFxUX1xcZVxcbSIsMCx7ImN1cnZlIjotNH1dLFsxLDNdLFszLDIsIlxcVF9cXGVcXGUiXSxbMSwyLCJcXFRfXFxlXFxtXlxcb3AiLDIseyJjdXJ2ZSI6NH1dXQ==
\begin{tikzcd}
{\T_\e\GL_\A} & {\T_\e\lnr_\A\times\T_\e\lnr_\A} & {\*} & {\T_\e\lnr_\A}
\arrow[from=1-1, to=1-2]
\arrow[from=1-2, to=1-3]
\arrow["{\T_\e\m}", curve={height=-24pt}, from=1-2, to=1-4]
\arrow["{\T_\e\m^\op}"', curve={height=24pt}, from=1-2, to=1-4]
\arrow["{\T_\e\e}", from=1-3, to=1-4]
\end{tikzcd}
\end{equation*}
In fact, $\T_\e$ is constructed by means of tangent pullbacks; however, as limits commute with other limits, tangent limits commute with other tangent limits. Furthermore, since $\lnr_\A$ is a differential object, $\T_\e\lnr_\A\cong\lnr_\A$. However, by Lemma~\ref{lemma:multiplication-and-unit-at-e}, $\T_\e\m=+$ and $\T_\e\e=0$. Moreover, since $+$ is commutative, it also follows that $\T_\e\m^\op=+^\op=+$. Therefore, we can rewrite the tangent triple equalizer diagram as the following tangent equalizer diagram:
\begin{equation*}
% https://q.uiver.app/#q=WzAsMyxbMCwwLCJcXGdsX1xcQSJdLFsxLDAsIlxcbG5yX1xcQVxcdGltZXNcXGxucl9cXEEiXSxbMywwLCJcXGxucl9cXEEiXSxbMCwxXSxbMSwyLCIrIiwwLHsib2Zmc2V0IjotMn1dLFsxLDIsIiEwIiwyLHsib2Zmc2V0IjoyfV1d
\begin{tikzcd}
{\gl_\A} & {\lnr_\A\times\lnr_\A} && {\lnr_\A}
\arrow[from=1-1, to=1-2]
\arrow["{+}", shift left=2, from=1-2, to=1-4]
\arrow["{!0}"', shift right=2, from=1-2, to=1-4]
\end{tikzcd}
\end{equation*}
However, by Definition~\ref{definition:maximal-differential-subring}, the tangent equalizer of $+$ and $!0$ is precisely the maximal differential subring $\lnr_\A^-$ of $\lnr_\A$. To conclude, we need to prove that the unique isomorphism $\gl_\A\cong\lnr_\A^-$ must be $\chi_\A$. This will prove that not only $\gl_\A$ and $\lnr_\A^-$ are isomorphic as differential objects but that they are also isomorphic as differential Lie algebras, since by Lemma~\ref{lemma:chi-vs-theta-maps}, $\chi$ is a Lie algebra homomorphism. To prove this statement, recall that $\chi\rho=\T_\e\theta$ as established by Lemma~\ref{lemma:chi-vs-theta-maps}, where $\<\rho,-\rho\>\colon\lnr_\A^-\to\lnr_\A\times\lnr_\A$ is the equalizer map of $+$ and $!0$. Let $\xi\colon\gl_\A\to\lnr_\A^-$ denote the unique isomorphism that satisfies $\xi\<\rho,-\rho\>=\T_\e\varphi=\<\T_\e\theta,\T_\e\i\T_\e\theta\>$. By postcomposing on both sides by $\pi_1$, we obtain that $\xi\rho=\T_\e\theta=\chi\rho$. However, since $\rho$ is monic, this implies that $\xi=\chi$. Therefore, $\chi$ becomes a differential Lie algebra isomorphism.
\end{proof}

When the differential object $\A$ admits negatives, then $\lnr_\A$ also admits negatives, with $-$ defined as the unique morphism that makes the following diagram
\begin{equation*}
% https://q.uiver.app/#q=WzAsNCxbMCwwLCJcXEFcXHRpbWVzXFxsbnJfXFxBIl0sWzEsMCwiXFxBIl0sWzAsMSwiXFxBXFx0aW1lc1xcbG5yX1xcQSJdLFsxLDEsIlxcQSJdLFsyLDMsIlxcZXYiLDJdLFszLDEsIi0iLDJdLFswLDEsIlxcZXYiXSxbMiwwLCJcXGlkXFx0aW1lcy0iXV0=
\begin{tikzcd}
{\A\times\lnr_\A} & \A \\
{\A\times\lnr_\A} & \A
\arrow["\ev", from=1-1, to=1-2]
\arrow["{\id\times-}", from=2-1, to=1-1]
\arrow["\ev"', from=2-1, to=2-2]
\arrow["{-}"', from=2-2, to=1-2]
\end{tikzcd}
\end{equation*}
commutative. Therefore, when $\A$ admits negatives, $\lnr_\A\cong\lnr_\A^-$. In this case, Theorem~\ref{theorem:main-theorem} becomes the classic result of differential and algebraic geometry that establishes that the Lie algebra of the general linear group $\GL_n(\Real)$ is isomorphic to the Lie algebra $\lnr_n(\Real)$ of linear endomorphisms of $\Real^n$.

In particular, thanks to~\cite[Lemma~2.2.1 and Corollary~2.2.5]{macadam:vector-bundles}, when the tangent structure comes with negatives, in that the tangent bundle $\T M\to M$ of each object admits a negation $n\colon\T M\to\T M$, every differential bundle also comes with negatives. So, in particular, each differential object admits negatives.

\begin{corollary}
\label{corollary:main}
If $\A$ admits negatives, the Lie algebra of the general linear Lie group $\GL_\A$ of $\A$ is isomorphic to the Lie algebra $\lnr_\A\cong\lnr_\A^-$. In particular, in a Cartesian tangent category with negatives, $\gl_\A\cong\lnr_\A$.
\end{corollary}



%__________________________________________________________________________
%__________________________________________________________________________